\chapter[Sermon 90. The Judge, and Last Judgment Is Described, with the Resurrection of the Dead.]{The Judge, and Last Judgment Is Described, with the Resurrection of the Dead.}
\label{sermon:090}
\markright{Sermon 90. The Judge, and Last Judgment Is Described, with the ...}

And I saw a great white throne, and he who sat on it had a face from which the earth and heaven fled away, and their place was no more found. And I saw the dead, great and small, standing before God, and the books were opened. Another book was opened, which is the book of life, and the dead were judged according to what was written in the books, according to their deeds. And the sea gave up the dead who were in it, and death and hell delivered up the dead who were in them, and they were judged every man according to his deeds. And death and hell were cast into the lake of fire. This is the second death, and whoever was not found written in the book of life was cast into the lake of fire.

\index{Antichrist}John had begun to speak of the universal and final judgment, around the end of the eleventh chapter. And he resumed the same to be finished in the nineteenth chapter, where we heard that Antichrist should be cast down from his throne and glory into hell. A question then arose concerning those who, although they do not cleave to Antichrist, are also not joined with Christ: what will become of them at the final judgment? Having addressed that question and demonstrated the equity of God’s judgments, he returns, as it were, with an after-song to the description of the general and final judgment, describing it more generally now than he did in the nineteenth chapter, where he seemed chiefly to have treated of the destruction of Antichrist; yet so that he also showed, in a manner of speaking, what would happen to the other ungodly.

Now he handles the same judgment more generally, showing that all shall be judged herein, and sets it forth wholly as if it were painted to be seen by our eyes. For in his customary manner, he explains all this matter through a heavenly vision, so that he might not seem only to tell the thing to our ears, but also to show it forth to be seen by our eyes, so that it might be more deeply imprinted in our minds. And all these things are most certain and undoubted—as I also admonished you before—revealed by the judge Christ himself. But the judge and Lord himself can be ignorant of nothing concerning this matter. Nor can we perceive that John has hitherto been deceived or misled in anything that he has set forth to us, but has hit rightly all and singular points, as we see that testimonies confirm his prophecies to be fulfilled. Why, then, should we so much as doubt one of such things as are spoken of the judgment?

\index{Resurrection}Therefore let us credit these things, and not be among the mockers whom the Apostle Peter prophesied would come and say: “Where is the promise of his coming?” Doubtless this matter is of greatest importance, the foundation and root of our faith. Here are explained to us not a few articles of our sincere and catholic faith, chiefly these: “I believe that Christ shall come to judge the living and the dead; I believe the communion of saints, the resurrection of the flesh, and life everlasting.” Let us therefore be diligent in hearing and marking these things, lest we be accounted among those who hear without any fruit the mysteries of the kingdom of God; but let us rather prepare ourselves to meet the judge, so that we may, with the wise virgins, enter with the bridegroom to the marriage and everlasting joys.

And the description or demonstration of this vision chiefly concerns: what the judge will be; who will be judged; how they will be judged; of what sort the resurrection of the dead will be; and of everlasting damnation; finally, who will be properly damned. These things I shall, in order and according to the grace that God has given me, declare as plainly as I can.

\index{Arethas of Caesarea}\index{Daniel (Book of)}We have previously understood what kind of judge there will be; at present, he is shadowed by certain marks. These things align with the vision described by Daniel in chapter 7. By the way, we see again how this book contains testimonies of the prophets, from whom it is commended to us, just as John explains the prophets to us. John sees a throne, and it is white and great. For the judge himself said that he would come in glory and majesty, namely with great light. And we also believe that his judgments are righteous, just, and white. Arethas, an interpreter, says that the throne is great because he who sits upon it is great, of whom the prophet said, “Great is the Lord, and great is his power,” etc. And in the throne, as judge of all and supremely righteous, he sits, furnished with all power and virtue. For all of this signifies the word “sitting.” Those who are to be judged stand; he sits. Therefore, he calls him who sits, as one would say, “judge.” No other name does he give. But we believe that all judgment is given to the Son, and that he is appointed judge over all. John therefore sees, and also shows us to behold the Lord Jesus Christ coming in the clouds of the air, a righteous and mighty judge. Paul also in 2 Timothy calls him a great God: not that there is one great God and another lesser God, but that the majesty of our Lord Jesus Christ will at that day be most evidently seen, and the Lord himself will then reveal himself to the world with greater glory and power than ever before.

\index{Apocalypse (Book of)}The same shall appear also most severe and most just. Whereupon John says figuratively, from whose face both heaven and earth fled away. For if those things which have not sinned dare not come into the judge’s sight, but seek as it were to save themselves by flight: where, I pray thee, shall the ungodly and summer [?]? And doubtless the prophet Malachi also: who, he says, shall abide the day of his coming? Or who is able to stand when he shall appear? So in the sixteenth chapter we heard that heaven fled back and was folded up like a scroll, that the mountains also and islands flitted, and that kings and princes and other men hid themselves in caves, saying to the hills and rocks, “Fall upon us, hide us from the face of him that sits on the seat, and from the wrath of the Lamb,” and so forth. By which words although a desperate conscience out of corrupt doctrine is described, yet the same shall appear chiefly in this judgment, when the severe and most righteous judge shall appear. A similar figure is read in the eighteenth Psalm. Where it is added, “and their place was no more found.” It is annexed to amplify the matter, not that heaven and earth shall be nowhere, but so much as they dare not (which is spoken by a figure) appear in the judgment of God. By all these things therefore is signified that the ungodly, being destitute of all counsel, shall not at that day know whither to turn themselves, or what to do, but trembling and despairing to be vexed with unspeakable torments before the seat. It might be thought in the meantime that John signifies this also, how heaven and earth should at the coming of the judge be renewed. Which also the Apostle Peter more plainly expresses in the third chapter of the second Epistle, which nevertheless refers and applies all those his sayings to the same sense that we have touched before. For he says: “Seeing then that all these things shall be dissolved, what ought you to be in holy conversation, looking for and hasting the coming of the day of God?” Arethas of Caesarea says that the flight of heaven and earth signifies no changing of place (for whither should they flee?), but flight and flitting from corruption to incorruption, and the last coming of the Lord, under which this mortal body of ours shall put on immortality, and the face of the earth shall be renewed. This, he says, is a like phrase of speech as is had in the twelfth of the Apocalypse, of the angels cast down out of heaven: neither was their place found any more in heaven.

Now he also touches those who will be judged, truly the dead. For he says, “And I saw the dead.” And shortly after, we shall hear that the dead will be raised up. Therefore, those who rise from the dead will be judged. Nevertheless, the living are not exempted, as the Apostle most manifestly states in 4.1 of 1 Thessalonians. But at this present, he names not the living, but the dead, for the resurrection of the dead is more hardly believed; and it is more easily believed that those remaining in flesh will be judged at that day. And truly, souls never die, bodies die. Therefore, where it is said here that the dead will be judged, we mean that all those who are dead at that day will come in their own bodies to the judgment of Christ. And all men must be judged. Wherefore, John sees great and small: that is to say, men of all sorts, status, sex, and age. Kings and princes are not exempted, the common people will not escape, neither children nor old folks, men nor women. He sees them all standing before the face, or judgment seat, of God. The guilty, or those accused, or to be accused, will be set before the judgment seat of God. And Paul also testifies expressly of this matter: “We must all appear before the judgment seat of Christ, that every one may receive in his body according to what he has done, whether it be good or evil.” 2 Corinthians 5. But they will appear in a diverse manner, both good and evil. For the wicked, as guilty, are brought to be judged and punished, and that their guiltiness may be openly known to all creatures. The good, for as much as they are justified and acquitted, and have no more guilt or crime by reason of Christ’s satisfaction, appear in judgment with glory, ready to judge the ungodly after their fashion and manner, and not to be judged by any. And this thing is singular: he says that we shall be judged in the sight of God. For who can appear in the sight of the terrible God, and fire consuming all things, save he that is purged with the blood of Christ? And what shall we think can be hid or escape the sight of God, knowing all things?

John also declares how the dead will be judged: he says that books are opened, and another book is opened, and so on. Therefore, by the books, and then by the book of life: that is to say, by the things written in those books, the dead are judged. Scripture ascribes to God the ways of humans, by which humans are accustomed to write down reminders for themselves, lest they forget things; but with God all things are one and always present; he neither forgets nor remembers. Nevertheless, Scripture attributes both to him. God is said to forget when he does not help or punish; again, he is said to remember when he does help or punish. In Malachi, the ungodly say that God has no care for human matters, neither does he care for the godly, nor yet punishes the wicked. But immediately an answer is made: then those who feared the Lord spoke to every one of his neighbor: the Lord gave ear and heard, and a book of remembrance was made in his presence, and so on. Therefore, their books opened, that is to say, the secrets of all people brought to light or made manifest, the Lord will judge whatever has been thought, said, done, or left undone. The books also of consciences (for the conscience stands in place of a thousand witnesses) will be opened in judgment, God revealing and judging all things. For Paul, speaking of the Gentiles, says that they show the work of the law written in their hearts, their conscience also bearing witness, and their thoughts accusing one another, or also excusing, in that day when the Lord will judge the secrets of men, according to my gospel, through Jesus Christ. And these are indeed the books which will be unclosed in the judgment. From this it appears that the judgment will be done with great speed; neither will every man be reasoned with by all, by books written to make the judge weary, as the ignorant might imagine.

But what is that unique book of life, which will also be opened at the judgment? Reference is made to the book of life in chapter 3. Briefly, the book of life contains only one entry: whoever believes in the Son of God has eternal life. Therefore, we are judged according to what is written in the book of life. Those who believe are saved; those who do not believe are already judged, that is to say, are certainly damned.

\index{Turks}\index{Judgment, Last}And for as much as faith shows itself by works, incredulity also, hidden in the heart, betrays itself by works: therefore John adds immediately, “according to their works.” For in Scripture, man is likened to a tree. And the tree is judged by its fruit, whether it be good or evil. A tree has a growing life, which in Latin is called \textit{anima vegetativa}, and a nature or disposition, bringing forth fruit according to its nature and kind. But that soul, that \textit{anima vegetativa}, and that good disposition, bringing forth in us good fruit, that is to say good works, is a lively faith in Christ, where it is, there the man is regenerated and has a good disposition: therefore he cannot, because of his good disposition, but bring forth good fruit. Therefore, after our works we shall all be judged. For the judgment must be open and manifest: but faith does not appear, but in works. For it is the gift of God, and is of itself invisible, namely, a sure trust in the promises of God. And it is seen in works. However, it does not follow from this that men are justified by works also, and not by faith alone: but that by works faith is declared, which purifies and justifies, that afterward we may be able to bring forth the works of righteousness. It follows that in judgment no pretense, no hypocrisy shall be allowed. For many say they believe, which declare their faith by no good works. We learn hereof that no book shall be of force at the last judgment, save the books of God, or the books of conscience, wherein God writes with his finger: finally, the book of life written by God before the worlds were made, through his divine predestination, whereby he has predestined us, that he might adopt us for his children by Christ Jesus. And the rest, which Paul recites in 1 Ephesians. Therefore, the hurtful books of Jews, Christians in name only, and Turks, such as the Talmud, decretals, and Alcoran, shall perish. These shall be of no force at all in the judgment.

Now he returns to the dead, of whom he had made mention before, and lest any man should say: how shall the dead be judged, who were drowned in the sea, who were swallowed up by fishes, and devoured by wild beasts, who were consumed with fire, or in the earth, were brought into dust? he anticipates and declares that the bodies of the dead will rise again, and being so restored will come to judgment, and says: and the sea gave up the dead that were therein; that is to say, those who had perished in the sea. And by these words also he has touched upon the manner and means of the resurrection of the dead, and has sent us also to Genesis 1. The manner of the resurrection is God’s omnipotence, as Saint Paul also witnesses in 3 to the Philippians. For God by his omnipotence raises up, and calls those things that are not, that they may be. If this thing seems to you new or impossible, behold the beginning of things, and from that estimate the small restitution. Was not the sea or water from the beginning? But is it written to have had any fishes from the beginning? None at all. But God commanded that the water should be replenished with fish. And did not immediately, at God’s command, all manner of fishes appear, where before there was not one? What marvel is it then, if God in the end of things, commands the sea, and other elements also, to yield again their dead, and they obey their maker? Verily the Lord in the gospel says that those who are in their graves also shall hear the voice or command of the Son of God, and shall rise again. The bodies moreover of those who die are turned for the most part into the same elements from whence they were taken out. There is that putrefies in the earth and is converted into earth. There are some consumed with fire. There are some that perish in water. Some hang in the air and are there consumed. But at the Lord’s command, by whatever kind of death they perish, they shall rise again to judgment whole. Arethas, also Bishop of Caesarea, perceived this and said: he recites these things to the intent he might declare what the final and universal resurrection shall be. For where many, not believing that the same shall be, say that it is by no means possible to be in those bodies which have been long corrupted and brought to that point that they are not at all, this sermon now corrects this, saying: Like as the bodies, when they were not, began to be, not by a certain chance or of themselves, but of the four elements, namely of Water, Fire, Air, and Earth: so also being reasonably returned again into the same, may be of the same composed again.

And for a further declaration he adds again: and death and hell gave up those who were in them, dead. For he understands by death, any kind of death, as though he should say: death itself restores to the Judge and judgment, whomever, after what sort soever he has dispatched. Death therefore is feigned to be as it were a person, which holds the dead in himself, or in a prison. And hell has yet but a few bodies (for some we read to have gone down to hell quick) but the souls of the wicked. The same return to their bodies, that the whole man may be judged, body and soul. Others by hell, after the Hebrew phrase, understand a sepulture or grave. Again is repeated, that the whole man shall be judged body and soul, after every man works.

\index{Primasius}Thus much hitherto concerning the resurrection of the dead, of which in our books elsewhere we have treated more at length. Finally, there follows a discussion of everlasting damnation, and who are properly condemned. And Hell, he says, and death are cast into the lake of fire. Of this, it has been spoken before. And Hell here signifies not the place of punishment, but those who are inhabitants of Hell, namely, whose souls are yet detained in hell, or appointed thither. Death also signifies those who are dead in sin, and those who, from spiritual or temporal death, go straightway to everlasting death. Whereupon is immediately annexed: This is the second death, by which truly those who are dead to Christ are devoted to perpetual fire, and those who live to Antichrist and the world. Others interpret these things to mean that after the judgment the Saints shall neither be buried again nor die. Paul affirms this also from Hosea in 1 Corinthians 15. Arethas and Primasius agree with us. For Arethas says: and he calls death and hell those who have committed things worthy of punishment, as fulfilling the number of the second death. And Primasius, by these names, says he, he signifies the Devil (because he is the author of death and pains in Hell) and also the whole fellowship of devils. For this is the same that he spoke more plainly before, by way of preventing: and the Devil, who deceived them, was cast into the lake of fire and brimstone. And that which he added there more obscurely, saying, and the beast and the false prophet, here is stated more plainly. So much Primasius. And who knows not that the members must follow the head, that all the ungodly follow the Devil, the head of all ungodliness?

And he expresses this most clearly, indicating who will be condemned to eternal fire at the final judgment: those whose names are not found in the book of life. Therefore, only those who are faithful in Christ, those predestined to eternal life, will be saved. All others, regardless of their religion or the kind of life they have lived, however righteous it may seem, will perish. Some interpret these words as referring to those who remain alive at that day. For we believe that the Son of God will judge both the living and the dead. Whether living or dead, it is certain that no one will be saved except through faith in Jesus Christ; all others will be damned. This is the ultimate fate of the good and the evil. To Jesus Christ, judge of all and redeemer of the faithful, be praise and glory forevermore. Amen.
